%!TeX root=../tese.tex
%("dica" para o editor de texto: este arquivo é parte de um documento maior)
% para saber mais: https://tex.stackexchange.com/q/78101

%% ------------------------------------------------------------------------- %%

% "\chapter" cria um capítulo com número e o coloca no sumário; "\chapter*"
% cria um capítulo sem número e não o coloca no sumário. Alguns autores
% preferem que a introdução não seja numerada, mas ainda assim apareça
% no sumário. Para isso, você pode usar "\chapter**{Introdução}". Esse
% comando não é padrão do LaTeX; ele foi definido neste modelo.
\chapter{Introdução}
\label{cap:introducao}

O desenvolvimento do kernel Linux constitui um exemplo relevante de produção colaborativa e distribuída de software livre (\cite{ferreira2019maintainers}). Nesse contexto, contribuidores, revisores e mantenedores colaboram por meio de duas infraestruturas distintas: as listas de \textit{e-mail}, nas quais os \textit{patches} são submetidos e revisados, e os repositórios Git, que registram as alterações aceitas e integradas ao código-fonte. Após eventuais modificações, as versões aprovadas são incorporadas às árvores dos mantenedores e avançam por uma cadeia hierárquica de integração até a \textit{mainline} (\cite{ferreira2019maintainers, ramsauer2019list}). Em razão dessa separação, nem sempre existem vínculos explícitos entre um \textit{e-mail} de submissão, as sucessivas versões do \textit{patch} e o \textit{commit} resultante (\cite{jiang2014tracing}).

Além de constituir uma representação técnica da alteração, o \textit{e-mail} de submissão exerce uma função comunicacional, pois apresenta a motivação, a solução proposta, os testes realizados e as eventuais limitações da modificação. A qualidade dessas informações influencia a revisão, a decisão de integração e a compreensão histórica do código (\cite{tan2019communicate}). As discussões associadas também registram justificativas, críticas e negociações entre desenvolvedores e mantenedores. Assim, problemas de comunicação podem interromper a revisão ou contribuir para a rejeição da proposta (\cite{ferreira2021incivility}).

A participação dos mantenedores amplia essa dimensão sociotécnica. Além de revisar o código, eles selecionam alterações, solicitam novas versões, acrescentam trailers de revisão e teste e definem a árvore em que o \textit{patch} será aplicado. Consequentemente, as listas documentam tanto o conteúdo técnico quanto as interações humanas que orientam a evolução das contribuições. Respostas, revisões sucessivas e trailers, como \textit{Reviewed-by}, \textit{Tested-by} e \textit{Signed-off-by}, podem, portanto, fornecer evidências adicionais sobre a trajetória de um \textit{patch} (\cite{ferreira2019maintainers}).

Embora a árvore Git registre os \textit{commits} integrados, ela não preserva necessariamente todo o histórico anterior à integração. O \textit{commit} representa o resultado consolidado do processo, enquanto as listas de discussão contêm os artefatos de sua evolução (\cite{ramsauer2019list}). Durante a revisão, os \textit{patches} podem ser corrigidos, reenviados, combinados, divididos ou adaptados às mudanças ocorridas na base de código. Além disso, o mantenedor pode modificar a mensagem, resolver conflitos ou alterar o código antes da integração. Operações de \textit{squash}, \textit{cherry-picking} e \textit{rebase} aumentam essa complexidade, permitindo que vários \textit{e-mails} originem um único \textit{commit} ou que uma proposta seja reorganizada em múltiplos commits. Dessa maneira, a relação pode apresentar cardinalidades de um para um, um para muitos ou muitos para um, e o vínculo entre o Message-ID e o hash do \textit{commit} pode ser perdido (\cite{ramsauer2019list}).

Nesse cenário, a reconstrução da trajetória de um \textit{patch} configura um problema de vinculação de registros, pois as versões submetidas às listas e o \textit{commit} integrado podem apresentar diferenças estruturais e textuais. A comparação integral dos artefatos é insuficiente: alterações equivalentes podem conter índices Git e linhas de contexto distintos, enquanto versões sucessivas podem modificar conjuntos apenas parcialmente coincidentes. \cite{ramsauer2019list} representam cada \textit{patch} pela combinação entre a mensagem e o diff, comparando separadamente os componentes textuais e as modificações efetuadas. \cite{jiang2014tracing}, por sua vez, demonstram que a comparação das linhas adicionadas e removidas apresenta desempenho adequado na vinculação entre \textit{e-mails} e \textit{commits}, enquanto os \textit{checksums} contribuem para relacionar versões sucessivas de uma mesma submissão.

A análise, entretanto, não deve se restringir às semelhanças entre os diffs. Versões relacionadas podem estar distribuídas por diferentes \textit{threads}, formando históricos cuja recuperação permite identificar alterações aceitas ou rejeitadas e examinar as interações que influenciaram sua evolução. \cite{jiang2014tracing} verificaram que aproximadamente 25\% dos \textit{patches} analisados apresentavam históricos de revisão superiores a quatro semanas e que as submissões com múltiplas versões tendiam a ser maiores e a apresentar maior taxa de aceitação. Ademais, a proximidade temporal entre \textit{e-mails} e \textit{commits} não comprova sua correspondência semântica. Por essa razão, devem ser considerados o conteúdo das mensagens, as conversações completas e as características técnicas das alterações (Bock et al., 2022).

A escalabilidade constitui uma dificuldade adicional, uma vez que a combinação irrestrita entre centenas de milhares de \textit{e-mails} e \textit{commits} produz um espaço de comparação computacionalmente impraticável. Para reduzi-lo, \textcite{ramsauer2019list} empregaram janelas temporais, arquivos modificados em comum e o agrupamento prévio de versões semelhantes. A parametrização avaliada para o kernel Linux alcançou índice de Fowlkes-Mallows de 0,911 em relação a uma referência construída manualmente e de 0,988 na validação com a lista \textit{linux-commits-tip}. Os autores também observaram que 99,5\% dos patches integrados chegaram ao repositório em até um ano, 80\% em até 40 dias e 50\% em até uma semana. Esses resultados fundamentam o uso de janelas temporais na geração de candidatos, embora a proximidade cronológica não deva ser interpretada isoladamente como evidência de correspondência.

A filtragem temporal deve, portanto, ser complementada por critérios estruturais, textuais e contextuais. Bock et al. (2022) demonstram que somente entre 29\% e 47\% das atividades de comunicação e codificação temporalmente próximas apresentam relação de conteúdo, o que reforça a necessidade de analisar conversações completas, e não apenas mensagens isoladas. De modo semelhante, \cite{ferreira2021incivility} combinaram assunto, autoria, ordem temporal, arquivos e proporção de linhas modificadas, obtendo precisão de 98,51\% e revocação de 90,08\% na identificação de associações entre revisões e \textit{commits}. Desse modo, uma estratégia escalável deve integrar filtragem temporal, similaridade do código e contexto comunicacional.

A necessidade de combinar essas evidências decorre da heterogeneidade das transformações sofridas pelos \textit{patches}. Quando o \textit{e-mail} e o \textit{commit} preservam os mesmos arquivos, linhas adicionadas e removidas e índices de blobs Git, a correspondência pode ser considerada praticamente determinística. Contudo, índices divergentes não implicam necessariamente alterações distintas, pois mudanças de base, \textit{rebases}, resoluções de conflitos e intervenções dos mantenedores podem modificar os conteúdos anteriores e posteriores. Assim, a coincidência dos índices constitui uma evidência forte, mas sua divergência não deve determinar automaticamente a rejeição de um vínculo.

Nos casos em que a correspondência não é exata, a decisão deve combinar evidências do código, dos metadados e do processo de revisão, como a similaridade das linhas alteradas, os arquivos em comum, o assunto e a mensagem, a autoria, os signatários, a sequência temporal, as referências a \textit{commits} corrigidos e os vínculos com as discussões. A correspondência deve, portanto, resultar da convergência ponderada de evidências técnicas, textuais, temporais, processuais e contextuais. Essa combinação permite reduzir falsos positivos sem excluir indevidamente relações verdadeiras que tenham sido modificadas durante a revisão e a integração.

Apesar dos avanços apresentados pela literatura, poucos estudos se dedicam a determinar, de forma automatizada e mensurável, se um \textit{e-mail} específico, publicado nas listas de desenvolvimento dos subsistemas do Linux, foi efetivamente integrado ao projeto. Dessa forma, este estudo busca preencher essa lacuna mediante uma abordagem capaz de verificar, com determinado grau de confiança, se um \textit{e-mail} contendo um patch foi posteriormente incorporado à mainline do Linux.
Para esse propósito, o trabalho utiliza dois conjuntos de dados complementares. O primeiro é o \textit{LKML5Ws: The What, When, Who, Where, and Why in the Linux Kernel Mailing Lists – A Columnar Dataset}, publicado no Zenodo por \cite{passos2026lkml5ws}. 

Esse conjunto reúne quase 30 milhões de \textit{e-mails} provenientes de aproximadamente 350 listas e abrange cerca de três décadas de desenvolvimento do kernel Linux. Sua organização colunar e seus arquivos anonimizados favorecem o processamento analítico em grande escala. Ao mesmo tempo, a distribuição das mensagens entre centenas de listas evidencia a complexidade da recuperação, da integração e da normalização desses dados (\cite{passos2026lkml5ws}). 
O segundo conjunto de dados é construído a partir de uma cópia completa do repositório principal do kernel Linux (\cite{torvalds_linux_source}), preservando o histórico de desenvolvimento disponível. Inicialmente, define-se uma revisão de referência, cujo identificador é registrado para assegurar a reprodutibilidade da coleta. Em seguida, extraem-se de cada \textit{commit} os identificadores, as relações de ancestralidade, a autoria, as datas, as mensagens, os arquivos modificados e as quantidades de linhas adicionadas ou removidas. Quando necessário, também são recuperados os diffs, os índices dos \textit{blobs} e os trailers de revisão.

Os dois conjuntos de dados representam, portanto, fases distintas e complementares do mesmo processo. O LKML5Ws registra a submissão e a discussão anteriores à integração, enquanto o conjunto de \textit{commits} documenta as alterações incorporadas ao histórico da \textit{mainline}. Com base nessa complementaridade, a abordagem proposta organiza a vinculação em duas etapas. Na primeira, são gerados pares candidatos mediante evidências de menor custo computacional, como janela temporal, assunto normalizado, arquivos em comum, compatibilidade de autoria, referências Link e índices Git. Na segunda, os candidatos são classificados independentemente por dois revisores. A consistência das classificações é examinada por meio da concordância Inter avaliadores, e as divergências são posteriormente resolvidas por consenso.

O resultado da vinculação não deve ser expresso apenas como uma decisão binária. Cada relação deverá receber uma categoria e um nível de confiança calibrado. Correspondências apoiadas por uma referência Link explícita e pela igualdade do assunto poderão receber confiança elevada. Casos sustentados pela igualdade dos índices Git, pela coincidência de autoria e por um intervalo temporal compatível poderão ser classificados como correspondências prováveis. Por outro lado, a inexistência de um \textit{commit} correspondente, mesmo após a aplicação dos critérios de busca, poderá indicar a rejeição ou a não integração do patch. Casos com evidências insuficientes ou contraditórias deverão ser classificados como inconclusivos, evitando que a ausência de correspondência seja interpretada automaticamente como rejeição.

Dessa forma, o objetivo geral deste trabalho é desenvolver e avaliar um método automatizado e mensurável para relacionar \textit{e-mails} que contenham \textit{patches} nas listas de desenvolvimento do kernel Linux aos \textit{commits} correspondentes incorporados à árvore principal. A contribuição pretendida é simultaneamente técnica e empírica. No âmbito técnico, propõe-se um procedimento reproduzível para geração, classificação e calibração de pares candidatos. No âmbito empírico, busca-se produzir evidências sobre o percurso das alterações entre sua submissão e integração. Ao restabelecer a ligação entre as listas e o repositório, a pesquisa procura recuperar parte da história pré-integração que se perde quando o \textit{patch} se converte em \textit{commit} e, assim, ampliar a compreensão sobre os processos de revisão e evolução do kernel Linux.

\textbf{Problema de pesquisa}

Apesar da disponibilidade pública dos arquivos das listas de desenvolvimento e da árvore Git principal, não existe uma ligação completa, uniforme e diretamente observável entre os \textit{e-mails} que submetem patches e os \textit{commits} resultantes de sua integração. O Message-ID identifica a mensagem, enquanto o \textit{hash} identifica o \textit{commit}; durante a revisão e a integração, contudo, esse vínculo pode desaparecer. Revisões sucessivas, alterações editoriais ou técnicas, reorganizações de séries, resoluções de conflitos, consolidações e divisões de \textit{patches} impedem que a relação seja inferida apenas por igualdade textual. Além disso, os conjuntos de dados disponíveis apresentam estruturas distintas, o que exige normalização, geração eficiente de candidatos, combinação de evidências heterogêneas e validação com base em uma referência confiável.

\textbf{Questão de pesquisa principal}

Como estabelecer, de forma automatizada e mensurável, a relação entre os \textit{e-mails} que contenham \textit{patches} nas listas de desenvolvimento do kernel Linux e os \textit{commits} correspondentes incorporados à árvore principal, atribuindo a cada vínculo um nível de confiança confiável e calibrado?

\textbf{Questões de pesquisa secundárias}


\begin{enumerate}[label=\textbf{\arabic*.}]
	\item Quais atributos dos conjuntos de dados LKML5Ws e dos \textit{commits} do kernel Linux apresentam maior poder discriminativo para distinguir correspondências verdadeiras de falsos positivos?
	\item Em que medida evidências potencialmente determinísticas, como referências Link, identificadores Message-ID e igualdade dos índices Git, abrangem as correspondências observadas?
	\item Qual estratégia de geração de candidatos reduz adequadamente o espaço de busca sem excluir uma proporção relevante de correspondências verdadeiras?
	\item Qual o ganho de desempenho do modelo calibrado em relação ao modelo treinado?
\end{enumerate}
